\chapter{Planetary pressures and conditional futures}\label{ch:planet}

\section{The valley has a boundary, but the world does not stop there}
Ridge and Vale are useful because we can draw them on one sheet. A water transfer appears as a decrease at one place and an increase at another. The real clinic, however, depends on conditions outside that drawing. Its electricity, equipment, medicines and food have histories elsewhere. Its water source depends on a larger landscape. Its waste crosses boundaries that the first diagram may not include.

Expanding the boundary does not mean that every local account must contain a complete model of Earth. It means that a claim must respect its boundary. A local transfer can be accurately described while remaining silent about upstream extraction or downstream pollution. The silence becomes misleading only when the local result is presented as a complete planetary judgement.

This chapter provides motivation for keeping those connections visible. It uses a small selection of dated scientific assessments, followed by original explanatory examples. The assessments describe pressures and possible futures. They do not validate EBU, identify its parameters or show that its proposed account can resolve the problems described.

\section{Read an environmental number with its date and unit}
A statement about planetary change should carry at least four pieces of information: what was measured or estimated, the period concerned, the comparison baseline and the relevant uncertainty or conditional assumptions. Without these, a number can travel into a claim it was never intended to support.

The IPCC's 2023 Synthesis Report assessed global surface temperature in 2011--2020 as $1.09\,^{\circ}\mathrm{C}$ higher than in 1850--1900, with a very likely range of $0.95$ to $1.20\,^{\circ}\mathrm{C}$. This is a global decadal comparison from that assessment. It is not the temperature of every location, a single-year value or a claim about the latest calendar year.\cite{intro-ipcc}

The distinction between a global average and a local experience is worth pausing over. A hypothetical average increase of one degree across two places could accompany different changes in each. The average is meaningful for its purpose, but cannot answer every question about local exposure. In the same way, an average resource stock does not reveal every empty tank.

The distinction between a decade and a year also matters. A long-period statistic smooths some short-term variation and answers a different question from a single observation. Choosing the appropriate timescale depends on the claim. A clinic manager planning tomorrow's delivery and a planner assessing infrastructure over decades need different kinds of information, even when both concern water.

\section{A projection is a conditional statement}
The same IPCC assessment gives different late-century warming estimates under different greenhouse-gas scenarios. For 2081--2100 relative to 1850--1900, its best estimates include $1.4\,^{\circ}\mathrm{C}$ for SSP1-1.9, $2.7\,^{\circ}\mathrm{C}$ for SSP2-4.5 and $4.4\,^{\circ}\mathrm{C}$ for SSP5-8.5. The corresponding very likely ranges are $1.0$--$1.8$, $2.1$--$3.5$ and $3.3$--$5.7\,^{\circ}\mathrm{C}$. These are conditional scenario results, not probabilities that societies will choose those pathways.\cite{intro-ipcc}

The word ``if'' is essential to this kind of reasoning. A conditional model says what follows under a specified collection of assumptions. It does not establish that those assumptions will hold. If a town asks how long a store would last under a declared withdrawal schedule, the answer changes when the schedule changes. That dependence does not make the calculation useless; it makes the assumptions part of the answer.

Nor is an uncertainty interval a permission to substitute any preferred conclusion. A range records a stated kind of uncertainty within an assessment. The uncertainty about a physical response and the uncertainty about future human decisions are not the same object. Combining them casually can make a conditional projection sound like a definite forecast or make a serious risk sound like an arbitrary guess.

For this book, the methodological lesson is more important than memorising a headline. Any future EBU study must likewise identify its forcing, initial condition, action rules and interpretation. A favourable result in one declared world would not become a forecast of every institution that uses the same acronym.

\begin{cautionbox}{Keep the condition attached}
``Under this scenario, the model projects this outcome'' is a different statement from ``this outcome will occur''. Removing the scenario converts a conditional result into an unsupported forecast. The same discipline applies to climate assessments and to the much smaller models developed in this project.
\end{cautionbox}

\section{Living systems cannot be reduced to one temperature number}
Climate is one dimension of planetary condition. The 2019 IPBES Global Assessment estimated that around one million animal and plant species were threatened with extinction. This was an assessment estimate of threatened species, not a census of species already extinct or a statement that every threatened species would disappear on a fixed date.\cite{intro-ipbes}

We do not need to turn this estimate into an EBU coordinate to recognise the modelling challenge it poses. A single homogeneous resource cannot represent the identities and relationships of different species. Exchanging one unit of an arbitrary scalar for another cannot by itself describe the loss of a particular ecological function. A future ecological application would require a much richer state and justification for any aggregation.

Consider an invented garden with two functions: food production and water infiltration. A change could improve one while worsening the other. Reporting only the first would give an incomplete picture; adding them requires a declared scale or decision rule. The example does not claim to model a real ecosystem. It shows why an apparently comprehensive score can hide choices about what counts and how.

The lesson for our introductory Gaussian model is therefore a boundary, not a defect to conceal. We use a conserved scalar to establish exact accounting relations. We do not announce that biodiversity is that scalar. The conceptual desire to preserve life-supporting conditions should make us more careful about representation, because those conditions include differences that simple addition can erase.

\section{Resources, services and material throughput}
The UNEP and International Resource Panel's \emph{Global Resources Outlook 2024} describes rising resource extraction and a conditional continuation in which extraction could increase by about $60\%$ between 2020 and 2060 without urgent concerted action. The headline is a global material-extraction comparison with a specified baseline and horizon. It is not a forecast of every country's consumption, nor an EBU study result.\cite{intro-resources}

An important conceptual question follows: what service is obtained from a material flow? Two hypothetical arrangements might provide the same number of refrigerated medicine-days while using different equipment lifetimes or maintenance practices. Conversely, equal material use could support different amounts or qualities of service. Throughput and function are connected but cannot be treated as synonyms.

A repair example makes the distinction concrete. Replacing an entire pump and replacing one failed seal may both restore flow under specified conditions. The material accounts differ. Whether the seal repair is actually adequate depends on inspection and performance, not merely on its smaller mass. A resource-efficient story is not evidence until the relevant function and duration are established.

This is also why ``use less'' is incomplete as a universal instruction. A place lacking safe infrastructure may require additional material investment to obtain essential service. Another arrangement may be able to reduce waste while preserving service. A responsible account needs to distinguish deprivation from avoidable throughput. The physical numbers inform that distinction; the purpose and institutional judgement must remain explicit.

\section{Three clocks in one decision}
The clinic operates on several timescales. Water must be available during today's procedures. A pump must remain serviceable through a maintenance interval. A watershed or energy infrastructure may be affected over much longer periods. A decision that looks favourable on one clock can impose a burden on another.

Suppose an invented temporary delivery solves today's shortage but requires repeated emergency journeys. The first journey may be necessary. The pattern may nevertheless suggest that a different arrangement deserves investigation. Calling the first action necessary does not make the long-term arrangement efficient; calling the arrangement inefficient does not make today's patient dispensable.

A present-state potential provides one declared view of the current arrangement. It does not automatically predict every delayed effect. To represent equipment wear, accumulated pollution or future availability, the state or an additional model must contain those quantities with suitable evidence. Merely enlarging the time interval in a sentence does not add them to an equation.

This distinction will matter when the book later proves a telescoping identity. A sequence of exact potential differences can be added exactly under stated conditions. The identity concerns the represented account. It is not a theorem that long-run services continue, that disturbances remain manageable or that every delayed ecological effect has been included.

\section{The importance of boundaries and displacement}
Imagine a town that reports a decrease in its own workshop waste after purchasing more finished components from outside. Local waste may indeed have fallen. Whether total waste associated with the town's use fell requires an account that includes the external production. Both statements can be investigated, but they have different boundaries.

This is an original teaching example of displacement, not a quantified claim about a particular trade pattern. It warns against evaluating a change through a boundary selected after its attractive result is seen. If the research question concerns a whole route, the route must be declared before the comparison. If it concerns one facility, the conclusion should remain at facility scope.

The same discipline applies to physical losses. A transfer cannot be called conservative merely because a source and destination are both listed. The sum must close within the declared coordinates or through an explicit boundary account. In the lossless teaching model, this is deliberately simple. In an extended model, a loss ledger would record where the represented quantity goes; it would not automatically become an additional term in the potential.

Nor does including a loss ledger justify counting the same effect twice. Conservation concerns the physical balance. Valuation concerns how the declared state is evaluated. Process burden concerns a residual declared effect not already represented in that evaluation. Keeping these roles separate will allow later chapters to make exact statements without pretending that every physical consequence is measured by one number.

\section{Why the future is neither guaranteed nor unknowable}
Environmental discussion can fall into two opposite errors. One treats a severe scenario as an inevitable destiny. The other treats uncertainty as a reason to avoid any conclusion. Conditional reasoning offers a more useful position: identify the plausible consequences of declared assumptions, examine uncertainty and ask which decisions remain sensible across relevant cases.

Our small teaching world offers an analogy. If Ridge can send water only through one pipe, a closed pipe prevents that transfer. We do not need to know the next year's exact demand to recognise the constraint. Yet knowing the constraint does not determine the best investment, which depends on alternatives, costs, priorities and evidence. Some facts are decisive at one level while leaving other questions open.

The EBU programme should be assessed with the same discipline. Exact accounting may be established before performance is known. Performance in a registered world may be observed before institutional usefulness is established. Neither uncertainty about the last question nor confidence about the first should be used to collapse the sequence.

The practical motivation is to improve the connection between decisions and the conditions they affect. That aim is compatible with uncertainty, provided uncertainty is recorded rather than hidden. It also demands that a proposal be open to revision if the representation or guidance proves unhelpful.

\section{A stock and a rate answer different environmental questions}
Imagine a hypothetical waste container with twelve units already inside it. A new process adds two units per day. The twelve is a stock; the two per day is a rate. Reducing the rate from two to one does not, by itself, reduce the existing stock. It slows the increase if nothing else changes. Removing three units per day would produce another account, which must also identify where the removed material goes.

The arithmetic is elementary, but the language can easily obscure it. ``Less added'' is not always ``less present''. ``Removed from this place'' is not automatically ``destroyed''. A complete account needs the initial stock, the flows and the boundary. This hypothetical example is not a model of a particular pollutant; different materials and processes require their own physical descriptions.

The same care will be needed in the EBU equations. A rate multiplied by a duration gives a quantity under the declared convention. Substituting a rate where a finite quantity is required changes the units and the result. Setting the duration to one can conceal the mistake numerically, which is why symbols and units should remain visible even in simple examples.

These distinctions also improve interpretation. A change that reduces an adverse inflow may be useful while leaving a large accumulated burden. A temporary improvement in a rate does not establish a completed recovery of the state. The account should preserve both facts instead of compressing them into one optimistic verb.

\section{Reading a headline as a careful student}
Suppose a headline states, ``Resource use will rise by sixty percent.'' A careful reader asks: which resources, relative to which year, by when, under what assumptions, and according to which source? The question is not an attempt to dismiss the finding. It restores the content needed to understand it.

Suppose another headline says, ``The model proves balance.'' The same habit applies. Which model? Which conserved quantity? Which definition of balance? Does the statement concern an algebraic identity, an equilibrium, a controller's performance or an empirical observation? The word ``proves'' should direct us to assumptions and an argument, not end the inquiry.

These habits help protect readers from both exaggerated criticism and exaggerated promises. A serious account can state severe risks without announcing that every detail is known. A serious research proposal can state exact mathematics without claiming to have solved those risks.

\section{Guided problems and answers}
\subsection{A temperature without a period}
A speaker quotes $1.09\,^{\circ}\mathrm{C}$ and calls it ``today's warming''. What is missing?

The figure used here concerns the IPCC's assessed 2011--2020 global average relative to 1850--1900. The period and baseline are essential. It should not be relabelled as a present-day annual observation. The answer illustrates why a correct number can still be used incorrectly.

\subsection{A conditional future}
A projection follows a declared high-emissions scenario. Does that establish that the scenario will occur?

No. The projection describes a response conditional on the scenario and the assessment's methods. It does not assign a probability to the human pathway merely by reporting a temperature range. Decisions can change the assumptions; physical uncertainty remains a separate issue.

\subsection{The clean local workshop}
A workshop reduces its local waste by buying a previously homemade component elsewhere. Has the whole route improved?

The local record alone cannot answer. The external component's relevant production and delivery must be considered under a declared boundary. The local improvement may be real while the wider claim remains unestablished. Expanding the claim requires expanding the evidence.
